% use the answers clause to get answers to print; otherwise leave it out.
\documentclass[12pts, answers]{exam}
%\documentclass[12pts]{exam}
\RequirePackage{amssymb, amsfonts, amsmath, latexsym, verbatim, xspace, setspace}

% By default LaTeX uses large margins.  This doesn't work well on exams; problems
% end up in the "middle" of the page, reducing the amount of space for students
% to work on them.
\usepackage[margin=1in]{geometry}
\usepackage{enumerate}
\usepackage{hyperref}
\usepackage{color,soul}

% Here's where you edit the Class, Exam, Date, etc.
\newcommand{\class}{NPRE 397}
\newcommand{\term}{Fall 2025}
\newcommand{\assignment}{Demonstrations of the Python for Reactor Kinetics Toolkit (PyRK)}
\newcommand{\duedate}{2025.12.10}
%\newcommand{\timelimit}{50 Minutes}

\newcommand{\nth}{n\ensuremath{^{\text{th}}} }
\newcommand{\ve}[1]{\ensuremath{\mathbf{#1}}}
\newcommand{\Macro}{\ensuremath{\Sigma}}
\newcommand{\vOmega}{\ensuremath{\hat{\Omega}}}

% For an exam, single spacing is most appropriate
\singlespacing
% \onehalfspacing
% \doublespacing

% For an exam, we generally want to turn off paragraph indentation
\parindent 0ex

%\unframedsolutions
\usepackage{bibentry}
\begin{document} 

% These commands set up the running header on the top of the exam pages
\pagestyle{head}
\firstpageheader{}{}{}
\runningheader{\class}{\assignment\ - Page \thepage\ of \numpages}{Due \duedate}
\runningheadrule

\class \hfill \term \\
\assignment \hfill Due \duedate\\ %\begin{flushright}
%\begin{tabular}{p{5in} r l}
%\end{tabular}
%\end{flushright}
\rule[1ex]{\textwidth}{.1pt}

%%%%%%%%%%%%%%%%%%%%%%%%%%%%%%%%%%%%%%%%%%%%%%%%%%%%%%%%%%%%%%%%%%%%%%%%%%%%%%%%%%%%%
%%%%%%%%%%%%%%%%%%%%%%%%%%%%%%%%%%%%%%%%%%%%%%%%%%%%%%%%%%%%%%%%%%%%%%%%%%%%%%%%%%%%%

For \class, your total grade will be earned through multiple computational 
exercises as well as contribution to a comprehensive suite of enhancements to 
the Python for Reactor Kinetics toolkit (PyRK). This independent study 
is intended to tie together ideas regarding reactor physics, reactor kinetics, data science, reproducible scientific 
compuing, and numerical methods for nuclear engineering.
\textbf{This project is more substantial and advanced in nature. It will 
require 6-10 hours of effort per week and will therefore earn 2 credit 
hours.}

The primary work will be to demonstrate the PyRK toolkit in the context of an 
advanced reactor design not yet covered by the suite of example problems. This 
effort may include resolving open issues in the PyRK repository and 
implementing improvements to the gaseous density model within. 
This work will emphasize usability and 
reproducibility of the toolkit, with special emphasis on demonstrating physics 
capabilities relevant to high temperature gas reactors (HTGRs). 
Software development can take place using personal or local
workstations, or if needed, cluster computing capabilities. Data and software 
contributions resulting from this work will be contributed to the open source 
repository at \href{github.com/pyne/pyrk}{github.com/pyne/pyrk}. 

The student will be expected to consistently handle assigned issues in the PyRK 
repository and to complete a final report and short presentation summarizing their accomplishments by 
the end of the semester. The project will be assessed as independent research work, much 
like a journal article undergoes peer review. I will be looking for : 

\begin{itemize}
\item Relevance
\item Technical Detail 
\item Analytic Rigor
\item Verifiability
\item Clarity
\item A Conclusion
\end{itemize}

This work will consist of four deliverables:

\begin{itemize}
        \item a work plan,
        \item simulation of reactor transients in an HTGR,
        \item a brief summary report in the style of a conference or journal paper,
        \item and a 10 minute summary presentation.
\end{itemize}

% ---------------------------------------------
\begin{questions}
\addpoints
% intro
\question[20] \textbf{Work Plan: Due 2025.09.05}

To help establish scope and milestones, the first step of the project will be a 
        work plan. This plan, developed under supervision of an Advanced 
        Reactors and Fuel Cycles Research Group Graduate Research 
        Assistant, will be delivered in the form of multiple GitHub 
        issues with descriptions and deadlines describing the work you intend 
        to do in the ARFC group. These issues should be compiled into a GitHub 
        Project in the ARFC organization.

Please arrange at least one short meeting with Prof. Huff to discuss the project plan
so that I can help you refine your schedule. I would be happy to provide feedback 
on a draft of your plan (once) before it is due.

\question[60] \textbf{Software Contributions: Due \duedate}
Complete the full suite of GitHub Issues agreed upon as a part of the Work 
        Plan, within the due dates specified in each issue. All issues must be 
        completed by \duedate. This will require responsive completion of all 
        necessary pull requests in  
        collaboration with colleages within the ARFC group. 

\question[20] \textbf{Final Report and Presentation: Due \duedate}
The final report must be submitted to Prof. Huff and a summary presentation 
        will be given to the ARFC group during a weekly group meeting late in 
        the semester. 

\end{questions}


%\bibliographystyle{plain}
%\bibliography{}

\end{document}
